\documentclass[11pt]{article}

\usepackage[margin=1in]{geometry}
\usepackage{amsmath}
\usepackage{amssymb}
\usepackage{hyperref}

\title{Stage 3I Physics Note: Small-\(x\) Regularity for Cartoon Variables}
\author{GL-with-AI Project}
\date{\today}

\begin{document}
\maketitle

\section{Scope}

Stage 3I documents regularity requirements near the cartoon/spherical axis
\(x=0\) for the future \(4+1\) black-string modified-cartoon implementation.
It is a documentation and symbolic-limit note only. It does not implement C++
source terms, finite differences, ghost-cell filling, or axis boundary
conditions.

\section{Conventions}

The physical spatial dimension is
\[
  \mathrm{GR\_SPACEDIM}=4,
\]
and the gridded Chombo dimension is
\[
  \mathrm{CH\_SPACEDIM}=2.
\]
The gridded directions are \(x,z\). The two hidden directions are equivalent
\(w\) directions, so
\[
  N_{\rm hidden}=4-2=2.
\]
The hidden component \(h_{ww}\) is Cartesian-like, not \(g_{\theta\theta}\).
The spherical/cartoon reconstruction supplies the external factor:
\[
  g_{\theta\theta}
  =
  x^2\gamma_{ww}
  =
  x^2\frac{h_{ww}}{\chi}.
\]
All traces are four-dimensional, and the conformal extrinsic-curvature
denominator remains \(1/4\).

\section{Cartoon Axis Versus Physical Singularity}

In the reduced black-string cartoon setup, the cartoon/spherical axis is at
\(x=0\). The unregularized black-string curvature singularity is also located
at \(x=0\). These two issues are co-located in the coordinate description but
are conceptually distinct.

The cartoon-axis problem is a removable coordinate-singularity, parity, and
tensor-regularity problem for smooth SO(3)-symmetric fields. The black-string
singularity is a physical curvature singularity of the target spacetime.
Turduckening or interior regularization is meant to remove or regularize the
physical singularity from the active evolved problem. The cartoon-axis
treatment separately handles the removable singular structures introduced by
the spherical/cartoon reduction.

Therefore turduckening addresses the fact that the physical singularity is
co-located with the cartoon axis, but it does not replace correct axis parity
and removable-singularity treatment. Conversely, passing the small-\(x\)
cartoon regularity checks does not prove that the black-string interior has
been physically regularized correctly.

\section{Apparent Singularities}

Cartoon source terms near \(x=0\) contain apparent singular structures such as
\[
  \frac{1}{x},\qquad
  \frac{1}{x^2},\qquad
  \frac{\gamma_{xx}-\gamma_{ww}}{x^2},\qquad
  \frac{\partial_x q}{x}.
\]
The corresponding conformal expressions involve \(h_{xx}\), \(h_{ww}\),
\(\chi\), and their derivatives. These expressions can be removable
singularities for smooth SO(3)-symmetric data, but they are unsafe in raw
numerical form unless the limiting behavior is enforced.

\section{Parity And Matching}

Scalar-like quantities such as \(\chi\), \(K\), \(\Theta\), and the lapse are
expected to be even in \(x\). Diagonal components \(h_{xx}\), \(h_{zz}\), and
\(h_{ww}\) are expected to be even. Components with one radial and one
non-radial index, such as \(h_{xz}\) and \(A_{xz}\), are expected to be odd.
The radial connection and shift components, such as \(\hat{\Gamma}^x\) and
\(\beta^x\), are expected to be odd, while the \(z\)-components are expected to
be even. These component parities should be confirmed against the chosen
GRChombo/Pau implementation conventions before coding.

The diagonal smooth-axis matching condition is
\[
  \gamma_{xx}-\gamma_{ww}=O(x^2).
\]
Since \(\gamma_{ij}=\chi^{-1}h_{ij}\), with \(\chi\) finite and even at the
axis, the corresponding conformal condition is
\[
  h_{xx}-h_{ww}=O(x^2).
\]
For the off-diagonal component,
\[
  \gamma_{xz}=O(x),\qquad
  h_{xz}=O(x),\qquad
  A_{xz}=O(x).
\]

\section{Hatted Connection Guard}

The highest-risk small-axis quantity is the hatted conformal connection,
especially \(\hat{\Gamma}^x\). Scalar-like variables and manifestly odd tensor
fields can be checked with generic Taylor representatives. In contrast,
\(\hat{\Gamma}^x\) is assembled from contracted angular Christoffels that can
carry bare \(1/x\) behavior, and the hidden-sector multiplicity
\(N_{\rm hidden}=2\) is load-bearing.

The local Stage 3I symbolic guard checks the documented singular cartoon
contribution
\[
  \tilde{\Gamma}^x_{\rm singular}
  =
  \frac{N_{\rm hidden}}{x}
  \left(1-\frac{h_{xx}}{h_{ww}}\right),
  \qquad
  N_{\rm hidden}=2.
\]
With \(h_{xx}-h_{ww}=O(x^2)\), this expression is finite and odd. If the
axis values of \(h_{xx}\) and \(h_{ww}\) do not match, the same assembled
expression exposes a residual \(1/x\) divergence.

The hatted connection convention is
\[
  \hat{\Gamma}^i
  =
  \tilde{\Gamma}^i+2Z^i.
\]
The symbolic check assumes regular \(Z^x=O(x)\) and verifies that
\(\hat{\Gamma}^x\) remains finite. The expected parity is that
\(\hat{\Gamma}^x\) is odd and \(\hat{\Gamma}^z\) is even, but the final sign
and full contracted-connection convention still require confirmation against
GRChombo's cartoon connection definition and Pau's implementation.

A residual axis error in \(\hat{\Gamma}^x\) can contaminate the Gamma-driver
shift/gauge sector, so this assembled check is stronger than generic
``odd-over-\(x\)'' tests.

\section{Symbolic Limits}

The Stage 3I SymPy check uses Taylor representatives
\[
  q(x,z)=q_0(z)+q_2(z)x^2+O(x^4),
\]
and
\[
  \gamma_{xz}(x,z)=s_1(z)x+s_3(z)x^3+O(x^5).
\]
It verifies limits such as
\[
  \lim_{x\to0}\frac{\partial_x q}{x}=2q_2(z),
\]
and
\[
  \lim_{x\to0}\frac{\gamma_{xx}-\gamma_{ww}}{x^2}
\]
is finite when the axis matching condition is imposed. It also verifies that
odd off-diagonal components have finite quotients such as
\[
  \lim_{x\to0}\frac{\gamma_{xz}}{x}=s_1(z).
\]
It also checks the assembled \(\tilde{\Gamma}^x\) singular part and
\(\hat{\Gamma}^x=\tilde{\Gamma}^x+2Z^x\) under regular \(Z^x\) parity. These
checks validate local removable-singularity bookkeeping only; they do not
derive the full CCZ4 RHS.

\section{Deferred Implementation Choices}

Stage 3I does not choose ghost-cell parity rules, one-sided stencils,
axis-point evaluation, turduckening cutoff handling, finite-difference order,
or C++ source-term organization. Those decisions are deferred to Stage 3J,
Stage 3K, and the later implementation stage.

\end{document}
